\documentclass{article}
\usepackage{amsmath,amssymb,amsthm}
% The nouns that the article uses and does not define, and the words that name numbers.
% Nothing is printed for them.
\usepackage[notions={graph, finite graph, connected graph, vertex (vertices)},
  quantities={degree, distance}]{leanify}
\newtheorem{theorem}{Theorem}
\newtheorem{lemma}[theorem]{Lemma}
\begin{document}

A test of notions that an article uses without defining them. The preamble names the
nouns \emph{graph} and \emph{vertex}, the adjectives \emph{finite} and \emph{connected} of
graphs, and the words \emph{degree} and \emph{distance} for numbers. A graph is then a kind of object of its
own, of which Lean knows nothing, and what is proved about graphs is proved whatever a
graph is.

\section{Quantities}

A sentence that introduces a symbol for a named quantity declares the quantity. Its name
says what kind of object it is: a number, or an object of a noun.

Let $\chi(G)$ denote the chromatic number of a graph $G$.
Let $\omega(G)$ denote the clique number of a graph $G$.
We write $\Delta(G)$ for the maximum degree of a graph $G$.
We denote by $d(G,H)$ the edit distance of graphs $G$ and $H$.
Let $L(G)$ denote the line graph of a graph $G$.

A name that does not say the kind is not read, and the finding says how to say it.

Let $\alpha(G)$ denote the independence ratio of a graph $G$.

\section{Results quoted from the literature}

A result whose heading cites its source and that the article does not prove is an
assumption of the article. The certificate names it among its axioms.

\begin{lemma}[{\cite[Section~5.1]{diestel}}]\label{clique}
For every finite graph $G$, $\omega(G)\le\chi(G)$.
\end{lemma}

\begin{lemma}[Greedy colouring, \cite{diestel}]\label{greedy}
For every finite graph $G$, $\chi(G)\le\Delta(G)+1$.
\end{lemma}

\section{What the article proves}

A step that uses only what it cites is checked, whatever a graph is.

\begin{theorem}\label{main}
Let $G$ be a finite connected graph. Then $\omega(G)\le\Delta(G)+1$.
\end{theorem}
\begin{proof}
By Lemma~\ref{clique}, $\omega(G)\le\chi(G)$. By Lemma~\ref{greedy}, $\chi(G)\le\Delta(G)+1$.
Hence $\omega(G)\le\Delta(G)+1$.
\end{proof}

A statement about an object of a declared noun, and about several.

\begin{lemma}\label{line}
For every graph $G$, $\chi(L(G))\ge\omega(L(G))-\omega(L(G))+\chi(L(G))$.
\end{lemma}

\begin{lemma}\label{distance}
For all graphs $G$ and $H$, $d(G,H)+1>d(G,H)$.
\end{lemma}

\section{What is not declared}

A noun that the preamble does not name is not read, and the check names it as one that
the article could declare.

\begin{lemma}\label{undeclared}
Let $T$ be a tree. Then $1\le2$.
\end{lemma}

\section{What is assumed of a kind of object}

A kind of object that the article does not define is assumed to have an element. A
sentence that says only that such an object exists is therefore not a claim, and the
finding says so. With a condition it is a claim like any other.

\begin{lemma}\label{exists}
There is a graph $G$. For every graph $G$ there is a vertex $v$.
There is a graph $G$ such that $\chi(G)\le\chi(G)$.
\end{lemma}

\begin{thebibliography}{1}
\bibitem{diestel} R. Diestel, \emph{Graph Theory}, 5th ed., Springer, 2017.
\end{thebibliography}
\end{document}
